\chapter{The shipping eraser and its specification (ErasureSpec, EraserAsks, the oracle reroute)}
\label{chap:eraser}

\section*{At a glance}
\begin{itemize}
\item \lead{Goal} Turn the eraser's oracle decision into a theorem; record every remaining
assumption about the world outside pure Lean terms.
\item \lead{Main results} Theorem~\ref{thm:isErasable-WF} (oracle sound at the built context);
Theorem~\ref{thm:Oracle-kernel_isErasable_sound} (a real run inherits it);
Theorem~\ref{thm:ErasureSpec-oracle_sound_of_run} (the shipping verdict is sound);
Theorem~\ref{thm:not_erasable_of_informative} (an informative inductive's value is never erasable).
\item \lead{Status} \stProved{} the oracle reroute; \stAssumed{} eight \code{ErasureSpec} fields
(class D), five \code{EraserAsks} fields (class C, two refuted in general); \stInherited{} through
lean4lean throughout.
\item \lead{Lean files} \code{Erasure.lean}, \code{Relevance.lean}, \code{RelevanceCheck.lean},
\code{CheckerAdequacy.lean}, \code{ErasureSpec.lean}, \code{CasesNames.lean}, \code{Origin.lean}.
\item \lead{Paper} [S, Sec.~7.3] for the eraser; [S, Sec.~7.4] for \S\ref{sec:eraser:origin}'s
case-step premises.
\end{itemize}

Earlier chapters relate terms; this one describes the \emph{program}, \code{Erasure.erase}, that
walks a \code{Lean.Expr} and asks an oracle whether each subterm is irrelevant. None of that is
inside Lean's logic, so this chapter states what is assumed, once, and turns the oracle question
into a theorem.

\section{The configuration, the state and the erasure monad}
\label{sec:eraser:config}
Five knobs configure the eraser; a predicate pins the one combination every theorem below assumes.

\begin{definition}[The erasure configuration]
\label{def:Erasure-ErasureConfig}
\lean{Erasure.ErasureConfig, Erasure.Config.Extern, Erasure.Config.Nat}
\leanok
\lead{In short} Five knobs select how the eraser compiles \code{Nat}, \code{@[extern]} constants
and three optional passes; defaults and pinned values in Definition~\ref{def:ConfigPinned}.
\begin{itemize}
\item \code{extern}: \code{.preferLogical} compiles from the Lean body; \code{.preferAxiom}
emits a body-less \lbox{} axiom for the target to realise;
\item \code{nat}: \code{.peano} compiles \code{Nat} as the unary inductive; \code{.machine}
rewrites it to 63-bit integers;
\item \code{csimp}: apply \code{@[csimp]} before erasure;
\item \code{remove\_irrel\_constr\_args}: prune irrelevant fields;
\item \code{auto\_inline\_typeclass\_dispatch}: mark typeclass artifacts for later inlining.
\end{itemize}
\end{definition}

\begin{definition}[The pinned configuration]
\label{def:ConfigPinned}
\uses{def:Erasure-ErasureConfig}
\lean{LeanToLambdaBox.ConfigPinned}
\leanok
\lead{In short} \code{ConfigPinned cfg} pins every field to what every correctness statement here
assumes.
\begin{center}
\begin{tabular}{lll}
\toprule
Field & Default & Pinned \\
\midrule
\code{extern} & \code{.preferAxiom} & \code{.preferLogical} \\
\code{nat} & \code{.machine} & \code{.peano} \\
\code{csimp} & \code{true} & \code{false} \\
\code{remove\_irrel\_constr\_args} & \code{false} & \code{false} \\
\code{auto\_inline\_typeclass\_dispatch} & \code{false} & \code{false} \\
\bottomrule
\end{tabular}
\end{center}
Three fields move away from default, so \hl{a bare \code{\#erase foo to "foo.ast"} runs a
configuration no theorem here is about}.
\lead{Lean} the binder \code{hcfg} of the capstone and of every green rung, discharged by five
reflexivities.
\end{definition}

\begin{remark}
\lead{Caveat} Machine-\code{Nat}, \code{@[csimp]} and \code{@[extern]}-axiom -- exactly what [R]
describes as implemented -- sit outside the verified perimeter by construction, not by omission.
\end{remark}

\begin{definition}[The erasure monad]
\label{def:Erasure-EraseM}
\uses{def:Erasure-ErasureConfig, def:Kername, def:GlobalDeclarations, def:InductiveId}
\lean{Erasure.EraseM, Erasure.ErasureState, Erasure.ErasureContext, Erasure.run, Erasure.liftMetaM}
\leanok
\lead{In short} \code{EraseM} is a state monad over a reader monad over \code{CoreM}: state
accumulates the output, reader carries the ambient scope.
\begin{description}
\item[State] the registered-inductives memo table with field masks; the Name-to-kername memo
table; the accumulating \lbox{} environment; kernames hinted for inlining.
\item[Reader] the local context; an optional map from each mutual-block member to its fresh
variable; the universe parameters, threaded to the oracle; the configuration.
\end{description}
\lead{Lean} \code{Erasure.run} starts from the empty state; \code{liftMetaM} runs a \code{MetaM}
action in the reader's context.
\end{definition}

\begin{remark}
\lead{Caveat} $\Box$ is \code{LBTerm}'s \code{Inhabited} default, so a \code{panic!} inside
\code{EraseM} does not abort: it prints to standard error, yields $\Box$, and the command still
exits successfully having written a wrong \code{.ast}. Sixteen sites do this.
\end{remark}

\section{The relevance oracle and its soundness}
\label{sec:eraser:oracle}
This section turns the oracle's per-node relevance decision from an assumption into a theorem.

\begin{definition}[The relevance criterion on the verified checker]
\label{def:isErasable}
\uses{def:lean4lean-grounding}
\lean{LeanToLambdaBox.isErasable, LeanToLambdaBox.isErasableProp, LeanToLambdaBox.isArityCheck, LeanToLambdaBox.isArityCheck.loop}
\leanok
\lead{In short} Erasable when the type is a proposition (a proof) or an arity (a type former) --
[S, Sec.~7.3]'s test, \hl{re-implemented on lean4lean's verified kernel} instead of Lean's
elaborator.
\begin{itemize}
\item \textbf{Proof branch}: infer the type via \code{inferType}, ask \code{isProp};
\item \textbf{Arity branch}: weak-head-normalise, peel one \code{forall} at a time, succeed at a
sort, fail elsewhere, \textbf{throw} on exhausted fuel so the caller decides, never a false
negative;
\item infer once, take the disjunction.
\end{itemize}
\lead{Lean} \code{isArityCheck} starts at fuel \code{ty.approxDepth + 1}.
\end{definition}

\begin{remark}
\lead{Caveat} Finding \textbf{F-DEPTH}: fuel (\code{approxDepth}, 8 bits, saturating at 256) can
exhaust behind an \code{@[irreducible]} alias, falling back and answering negatively -- a
completeness defect bounding \code{kernel\_ind\_head\_true}
(Theorem~\ref{thm:EraserAsks-oracle_informative}).
\end{remark}

\begin{definition}[The relevance decision the eraser makes]
\label{def:Erasure-isErasable}
\uses{def:isErasable, def:Erasure-EraseM}
\lean{Erasure.isErasable, Erasure.isErasableMeta}
\leanok
\lead{In short} The impure call the erasure family makes: run Definition~\ref{def:isErasable}
purely, at the ambient kernel environment, context and universe parameters.
\begin{itemize}
\item on success, return the verdict;
\item on a kernel error (unsupported construct, or exhausted fuel), \hl{fall back to the
elaborator check} \code{Meta.isProp} or \code{Meta.isTypeFormerType}, so the transpiler never
fails on that account.
\end{itemize}
\lead{Lean} \code{isErasableMeta} is the fallback.
\end{definition}

\begin{remark}
\lead{Caveat} Two same-named declarations: \code{LeanToLambdaBox.isErasable} is the verified pure
check (Definition~\ref{def:isErasable}); \code{Erasure.isErasable} is the impure call
(Definition~\ref{def:Erasure-isErasable}). The kernel arm's soundness is a \emph{theorem}
(Theorem~\ref{thm:Oracle-kernel_isErasable_sound}); the fallback's is an \emph{assumption}
(\code{oracle\_meta}), measured dead over $139{,}196$ constants, though that covers only the error
route.
\end{remark}

\begin{lemma}[Oracle soundness, proof disjunct]
\label{lem:isErasableProp-WF}
\uses{def:isErasable, def:Erasable, def:lean4lean-grounding}
\lean{LeanToLambdaBox.isErasableProp.WF}
\leanok
\lead{In short} If \code{e} translates to \code{e'} and the \hl{proof-branch check succeeds on
\code{e}}, then \code{e'} is \Erasable{} by its proof disjunct: its inferred type lives in
\code{Sort 0}.
\lead{Lean} a Hoare-style obligation over an arbitrary well-formed checker state, not one run.
\end{lemma}
\begin{proof}
\leanok
\uses{asm:lean4lean-trust, asm:lean-reflection-axioms}
\lead{Idea} Compose lean4lean's well-formedness lemmas for \code{inferType} and \code{isProp}: a
positive verdict gives a typing at \code{Sort 0}, the proof disjunct of \Erasable{}.
\end{proof}

\begin{lemma}[Arity-up-to-definitional-equality is closed under forall-introduction]
\label{lem:IsArityUpTo-forallE}
\uses{def:IsArity, def:lean4lean-grounding}
\lean{LeanToLambdaBox.IsArityUpTo.forallE}
\leanok
\lead{In short} If \code{B} is an arity up to definitional equality in the extended context and
both \code{A}, \code{B} are types, then \code{forall (\_ : A), B} is one too.
\lead{Lean} carries the environment's well-formedness and a context of types as ambient
hypotheses. \hl{The one metatheorem the arity branch needs that lean4lean lacks}, letting the
telescope recursion below reassemble a whole-type witness from the body's.
\end{lemma}
\begin{proof}
\leanok
\uses{asm:lean4lean-trust}
\lead{Idea} Untyped \code{forall}-congruence lifts \code{B}'s equality onto \code{forall A, B}.
\end{proof}

\begin{lemma}[Oracle soundness, type-former disjunct]
\label{lem:isArityCheck-WF}
\uses{def:isErasable, def:IsArity, def:lean4lean-grounding}
\lean{LeanToLambdaBox.isArityCheck.loop.WF, LeanToLambdaBox.isArityCheck.WF}
\leanok
\lead{In short} If \code{ty} translates to \code{ty'} and the arity loop succeeds at any fuel,
\code{ty'} is an arity up to definitional equality; likewise for the entry point, which starts the
loop at \code{ty.approxDepth + 1}.
\lead{Lean} both quantify universally over the fuel: neither claims it adequate.
\end{lemma}
\begin{proof}
\leanok
\uses{lem:IsArityUpTo-forallE, asm:lean4lean-trust, asm:lean-reflection-axioms}
\lead{Idea} Induction on the fuel.
\begin{enumerate}
\item Weak-head normalise; a sort leaf gives the base arity.
\item A \code{forall} node opens the binder, applies the hypothesis under it, and reassembles the
telescope witness with Lemma~\ref{lem:IsArityUpTo-forallE}.
\item Transport the witness back across the reduction.
\item Exhaustion throws, so the obligation is vacuous there.
\end{enumerate}
\end{proof}

\begin{theorem}[Relevance-oracle soundness]
\label{thm:isErasable-WF}
\uses{def:isErasable, def:Erasable, def:lean4lean-grounding}
\lean{LeanToLambdaBox.isErasable.WF}
\leanok
\lead{In short} If \code{e} translates to \code{e'} and the executable oracle succeeds, then
\code{e'} is \Erasable{}.
\lead{Lean} an obligation at an arbitrary well-formed state, applying to every run; its closing
note adds \hl{no new axiom, \code{sorry} or \code{native\_decide}}.
\lead{Status} \stProved{}, \stInherited{} through lean4lean.
\end{theorem}
\begin{proof}
\leanok
\uses{lem:isErasableProp-WF, lem:isArityCheck-WF, asm:lean4lean-trust, asm:lean-reflection-axioms}
\lead{Idea} Infer the type once and branch: \code{isProp} gives the proof disjunct
(Lemma~\ref{lem:isErasableProp-WF}), else the type-former disjunct
(Lemma~\ref{lem:isArityCheck-WF}).
\end{proof}

\begin{definition}[A verification context at an arbitrary ambient local context]
\label{def:Lean4Lean-TypeChecker-VContext-ofMLCtx}
\uses{def:lean4lean-grounding}
\lean{Lean4Lean.TypeChecker.VContext.ofMLCtx, Lean4Lean.TypeChecker.kernelNGen, Lean4Lean.TypeChecker.VContext.ofMLCtx_venv, Lean4Lean.TypeChecker.VContext.ofMLCtx_lparams, Lean4Lean.TypeChecker.VContext.ofMLCtx_vlctx}
\leanok
\lead{In short} \code{VContext.ofMLCtx} builds a lean4lean context at an \hl{arbitrary ambient
modelled local context}, generalising lean4lean's constructor, which hardwires the empty one.
\begin{itemize}
\item \code{kernelNGen} (prefix \code{\_kernel\_fresh}, index zero) \emph{reserves} every free
variable not of that shape -- every \code{\_uniq}-named identifier a real run produces;
\item all other fields are taken over unchanged.
\end{itemize}
\lead{Lean} three simp projections record the built context's environment, parameters and local
context as the given ones.
\end{definition}

\begin{remark}
\lead{Caveat} These five names, and the two of
Lemma~\ref{lem:Lean4Lean-TypeChecker-M-WF-run-prime}, sit in the \emph{dependency's} namespace
\code{Lean4Lean.TypeChecker} -- a documented deferral, and a hazard: a future lean4lean release
declaring \code{M.WF.run'} would clash. The reservation premise is \emph{cheap}, not vacuous.
\end{remark}

\begin{lemma}[Run adequacy at the ambient local context]
\label{lem:Lean4Lean-TypeChecker-M-WF-run-prime}
\uses{def:Lean4Lean-TypeChecker-VContext-ofMLCtx, def:lean4lean-grounding}
\lean{Lean4Lean.TypeChecker.M.WF.run', Lean4Lean.TypeChecker.VState.WF.initial}
\leanok
\lead{Given} every ambient free variable is reserved by the kernel's initial generator;
\lead{Then} the initial checker state is well formed at the built context, so a successful obliged
run there returns a value satisfying the postcondition.
\lead{Lean} \hl{lean4lean's own run lemma, generalised from the empty context to an arbitrary
one}.
\end{lemma}
\begin{proof}
\leanok
\uses{asm:lean-reflection-axioms}
\lead{Idea} Replace the empty context in lean4lean's own lemmas by the ambient one, taking the
extension witness at a reflexive lift with an empty equivalence manager.
\end{proof}

\begin{theorem}[The kernel arm of the oracle is sound]
\label{thm:Oracle-kernel_isErasable_sound}
\uses{def:isErasable, def:Lean4Lean-TypeChecker-VContext-ofMLCtx, def:Erasable, def:lean4lean-grounding}
\lean{LeanToLambdaBox.Oracle.kernel_isErasable_sound}
\leanok
\lead{Given} the environment model, the modelled context and the reservation premise are well
formed;
\lead{Then} a \hl{pure run} of the verified check, succeeding on \code{e} translating to
\code{ve}, witnesses \code{ve} is \Erasable{}.
\lead{Lean} phrases the run as an equation on the runner's result, not an obligation.
\lead{Status} \stProved{}, \stInherited{} through lean4lean.
\end{theorem}
\begin{proof}
\leanok
\uses{thm:isErasable-WF, lem:Lean4Lean-TypeChecker-M-WF-run-prime, asm:lean4lean-trust, asm:lean-reflection-axioms}
\lead{Idea} A three-line composition.
\begin{enumerate}
\item Theorem~\ref{thm:isErasable-WF} gives soundness at the built context.
\item lean4lean's well-formedness lemma for the runner discharges the fuel layer.
\item Lemma~\ref{lem:Lean4Lean-TypeChecker-M-WF-run-prime} converts it to a statement about the run.
\end{enumerate}
\end{proof}

\section{Preparation, registration, binders}
\label{sec:eraser:prepare}
A normalisation pass, then registration of inductives and binders into the emitted environment.

\begin{definition}[Pre-erasure preparation]
\label{def:Erasure-prepare_erasure}
\uses{def:Erasure-EraseM, def:Erasure-ErasureConfig}
\lean{Erasure.prepare_erasure, Erasure.replaceUnsafeRecNames, Erasure.remove_unsafe_rec, Erasure.name_occurs}
\leanok
\lead{In short} A normalisation pass, lifted from Lean's compiler frontend, run before erasure.
\begin{enumerate}
\item replace every \code{\_unsafe\_rec} constant by its safe twin;
\item honour \code{@[macro\_inline]};
\item inline auxiliary matchers;
\item honour \code{@[macro\_inline]} again -- matchers can expose strict \code{ite}/\code{dite};
\item only when \code{csimp} is on, \hl{apply \code{@[csimp]} replacements everywhere}.
\end{enumerate}
\lead{Lean} \code{name\_occurs} is the only recursiveness test, treating \code{forall}, sorts,
literals and variables as non-occurrences.
\end{definition}

\begin{remark}
\lead{Deviation} No MetaCoq counterpart; its doc comment warns it may ill-type the term under
\code{@[csimp]}, which \code{ConfigPinned} pins off. Finding \textbf{F-UNSAFEREC}:
\code{remove\_unsafe\_rec} is not injective, so a legal \code{mutual unsafe def} block can collapse
to a repeated key -- what \code{block\_keys\_distinct} excludes by assumption.
\end{remark}

\begin{definition}[The prepared passes]
\label{def:preparePasses}
\uses{def:Erasure-prepare_erasure}
\lean{LeanToLambdaBox.preparePasses}
\leanok
\lead{In short} The three \code{CoreM} passes: \code{replaceUnsafeRecNames}, Lean's
\code{macroInline}, Lean's \code{inlineMatchers}.
\lead{Lean} \hl{an index set, not the call sequence}: the eraser makes four calls, running
\code{macroInline} twice; membership covers the repetition, so three here and four in the code are
both correct.
\end{definition}

\begin{definition}[Registering an inductive block]
\label{def:Erasure-register_inductive}
\uses{def:Erasure-EraseM, def:Erasure-isErasable, def:Kername, def:InductiveId, def:OneInductiveBody}
\lean{Erasure.register_inductive, Erasure.addAxiom, Erasure.ConstructorArgMask, Erasure.InductiveArgMasks, Erasure.ConstructorArgRelevance, Erasure.filter}
\leanok
\lead{In short} Memoised registration of a Lean inductive with its mutual companions; on a miss:
\begin{enumerate}
\item synthesise the block's kername from its members' names, concatenated;
\item per constructor: a body-less axiom for \code{@[extern]}, and an argument mask -- all kept
unless \code{remove\_irrel\_constr\_args} is on, when \hl{each field is tested with the relevance
oracle};
\item dummy projection names for a single-constructor, single-member, non-recursive block;
\item cons one inductive declaration, carrying the block's parameter count.
\end{enumerate}
\end{definition}

\begin{remark}
\lead{Caveat} Finding \textbf{F-PROP}: every emitted inductive leaves \code{propositional},
\code{kelim}, \code{finite} at their defaults, so match-on-box never fires here. Unregistered:
\code{Erasure.filter} drops trailing elements, so a short mask silently truncates.
\end{remark}

\begin{definition}[Emitting binders]
\label{def:Erasure-mkLambda}
\uses{def:Erasure-EraseM, def:LBTerm, def:toBvar, def:BinderName, def:FixDef}
\lean{Erasure.mkLambda, Erasure.mkLetIn, Erasure.mkAlt, Erasure.mkDef, Erasure.fvar_to_name, Erasure.withLocalDecl, Erasure.withLocalDef, Erasure.lambdaMonocular, Erasure.lambdaOrIntroToArity}
\leanok
\lead{In short} Locally nameless: descend under a binder by minting a fresh variable, close it by
naming and abstracting back to a de Bruijn index.
\begin{itemize}
\item \code{withLocalDecl}/\code{withLocalDef}: the descent;
\item \code{fvar\_to\_name}: \hl{downgrades a binder name to anonymous unless every character is
ASCII-graphic} (the parser rejects other names);
\item \code{mkLambda}/\code{mkLetIn}: close one binder;
\item \code{mkAlt}: close a case alternative, binding the \emph{last} field to index zero;
\item \code{mkDef}: close one fixpoint member, principal-index zero.
\end{itemize}
\end{definition}

\begin{remark}
\lead{Caveat} Source of Chapter~\ref{chap:lowering}'s printable-binder clause; the stronger
all-ASCII form was \emph{refuted} at five rungs (Chapter~\ref{chap:capstone}). Two doc comments
still hedge on binder order; Chapters~\ref{chap:lowering},~\ref{chap:refinement} settle it, but
the hedge remains.
\end{remark}

\section{The erasure family and the entry point}
\label{sec:eraser:visit}
Eighteen mutually recursive functions, closed over one entry point.

\begin{definition}[The erasure family]
\label{def:Erasure-visitExpr}
\uses{def:Erasure-EraseM, def:Erasure-isErasable, def:Erasure-register_inductive, def:Erasure-mkLambda, def:LBTerm}
\lean{Erasure.visitExpr, Erasure.visitLiteral, Erasure.visitLambda, Erasure.visitLet, Erasure.visitProj, Erasure.visitApp, Erasure.visitConst, Erasure.visitConstApp, Erasure.visitCasesEta, Erasure.visitCasesEtaGo, Erasure.visitCtorEta, Erasure.visitCtorEtaGo, Erasure.visitConstructor, Erasure.visitAppArgs, Erasure.visitCases, Erasure.visitAlt}
\leanok
\lead{In short} One mutual block of sixteen \code{partial\_fixpoint} definitions (eighteen with
Definition~\ref{def:Erasure-visitMutual}'s two); \code{visitExpr} box-tests \emph{every} node
first: on success, $\Box$, with \hl{no recursion at all}. The table below covers all eighteen.
\begin{center}
\begin{tabular}{ll}
\toprule
Function & Role \\
\midrule
\code{visitExpr} & box test, else dispatch on head \\
\code{visitLiteral} & \code{Nat.succ} chain or machine int \\
\code{visitLambda} & emit a lambda binder \\
\code{visitLet} & emit a let binder \\
\code{visitProj} & re-index past pruned fields \\
\code{visitApp} & dispatch an application \\
\code{visitConst} & a bare constant, as zero-arg \\
\code{visitConstApp} & route to eliminator, constructor or plain \\
\code{visitCasesEta} & eta-expand an eliminator \\
\code{visitCasesEtaGo} & eliminator eta-expansion loop \\
\code{visitCtorEta} & eta-expand a constructor \\
\code{visitCtorEtaGo} & constructor eta-expansion loop \\
\code{visitConstructor} & emit a constructor \emph{applied} \\
\code{visitAppArgs} & erase an argument list \\
\code{visitCases} & emit an eliminator's case node \\
\code{visitAlt} & emit one case alternative \\
\code{visitMutual} & close the environment over one block
(Definition~\ref{def:Erasure-visitMutual}) \\
\code{get\_constant\_kername} & memoised kername lookup \\
\bottomrule
\end{tabular}
\end{center}
A free variable is emitted as an \lbox{} free variable, to be abstracted later. Sorts, products,
metavariables, de Bruijn variables: all unreachable. Metadata is skipped.
\end{definition}

\begin{remark}
\lead{Why} \code{partial\_fixpoint} (replacing an opaque \code{partial def}) yields equational
lemmas and the eighteen-motive \code{mutual\_fixpoint\_induct} that
Chapters~\ref{chap:refinement},~\ref{chap:coldstart} run on. Cost: nine proof-only monotonicity
lemmas and two behaviour-preserving rewrites, one leaving dead code behind.
\end{remark}

\begin{remark}
\lead{Deviation} The constructor case emits \emph{applied} form, never a block -- opposite
agda2lambox and MetaRocq's \code{constructors\_as\_blocks} -- so Chapter~\ref{chap:lambdabox}'s
semantics is pinned with the block-constructor flag off. [R]'s Listing~2 shows block form; \hl{the
shipping code disagrees, and the code wins}. Findings \textbf{F-SPARSE}, \textbf{F-ETA2} sit here
(table, \S\ref{sec:eraser:caveats}).
\end{remark}

\begin{definition}[The environment closure]
\label{def:Erasure-visitMutual}
\uses{def:Erasure-visitExpr, def:Erasure-prepare_erasure, def:Erasure-mkLambda, def:Erasure-EraseM, def:Kername, def:GlobalDeclarations}
\lean{Erasure.visitMutual, Erasure.get_constant_kername}
\leanok
\lead{In short} The remaining two members close the environment over what the subject reaches:
memoised kername lookup, erasing a mutual block on a miss, via the \emph{original} recursive
definition where one exists.
\begin{itemize}
\item no compiler value, or \code{@[extern]} under the axiom policy: a body-less axiom;
\item a lone declaration not mentioning its own name: erase it, cons a constant declaration;
\item otherwise: mint one variable per member, erase each body under that map, and \hl{register
each under its own kername with a bare, unapplied fixpoint block} at its index.
\end{itemize}
\end{definition}

\begin{remark}
\lead{Deviation} MetaCoq's \code{erase\_global\_decls}, on demand. Findings \textbf{F-ETA} (bare,
unapplied fixpoint bodies), \textbf{F-UNSAFEREC} (the stripped-name key map) and
\textbf{F-QUOT}/\textbf{F-EQREC} (the body-less axiom arm) sit here (table,
\S\ref{sec:eraser:caveats}).
\end{remark}

\begin{definition}[The entry point]
\label{def:Erasure-erase}
\uses{def:Erasure-visitExpr, def:Erasure-prepare_erasure, def:Erasure-EraseM, def:Erasure-ErasureConfig, def:ASTType}
\lean{Erasure.erase}
\leanok
\lead{In short} Prepare the subject, erase it from the empty state at the given configuration, and
return the untyped \lbox{} program with the kernames hinted for inlining.
\lead{Lean} \hl{the function Chapter~\ref{chap:capstone}'s capstone is about}.
\end{definition}

\begin{remark}
\lead{Caveat} \code{\#erase} wraps this and writes the \code{.ast} with the unverified
\code{LBTerm.to\_sexpr}: \hl{the theorem stack stops at the \code{Program} value}, not the bytes on
disk. The default configuration is not the pinned one (Definition~\ref{def:ConfigPinned}).
\end{remark}

\section{What the correctness statement assumes: ErasureSpec}
\label{sec:eraser:spec}
\code{ErasureSpec} names, once, every fact assumed about opaque Lean primitives.

\begin{definition}[Soundness of a relevance verdict]
\label{def:Oracle-MetaSound}
\uses{def:Lean4Lean-TypeChecker-VContext-ofMLCtx, def:Erasable, def:lean4lean-grounding}
\lean{LeanToLambdaBox.Oracle.MetaSound}
\leanok
\lead{In short} A positive verdict is sound at the ambient scope and at \emph{any} local context
agreeing with the given one and kernel-reserved: every lean4lean translation there is \Erasable{}.
\lead{Lean} \hl{what the fallback is assumed to deliver and the kernel arm proves} -- they meet in
Theorem~\ref{thm:ErasureSpec-oracle_sound_of_run}.
\end{definition}

\begin{definition}[The environment queries report what the environment holds]
\label{def:LookupAdequate}
\uses{def:isCasesOnName, def:Erasure-prepare_erasure}
\lean{LeanToLambdaBox.LookupAdequate, LeanToLambdaBox.CasesInfoAgreesK, LeanToLambdaBox.altNumFields}
\leanok
\lead{In short} The eraser's four environment queries (a constant's info, the compiler's
declaration info, a constructor's arity, an eliminator's metadata) report what the environment
holds, only advancing the fresh-name generator.
\begin{itemize}
\item declaration-info's block-membership is \hl{guarded by the name not being
\code{\_unsafe\_rec}} -- at such a name the answered block omits it;
\item the arity and metadata clauses hold both directions, excluding non-constructor,
non-eliminator heads.
\end{itemize}
\lead{Lean} \code{CasesInfoAgreesK}: \code{casesOn}'s metadata agrees with the declared block --
discriminant after parameters/motive/indices; one minor premise and one alternative per
constructor, each binding exactly its fields.
\end{definition}

\begin{remark}
\lead{Why} \code{Lean.Environment} has a private constructor: no Lean term denotes one, so no
statement about \code{getConstInfo} is provable in Lean. Hence an assumption, not a lemma.
\end{remark}

\begin{definition}[The kernel's blocks in the model]
\label{def:BlockAdequate}
\uses{def:Erasure-register_inductive, def:isCasesOnName, def:IndInfo, def:IndArity, def:CtorOf, def:ConstOrigin, def:CasesOnShape}
\lean{LeanToLambdaBox.BlockAdequate, LeanToLambdaBox.KernelFields}
\leanok
\lead{In short} The kernel's blocks, constructors and eliminators are the model's, at the \lbox{}
identifier the eraser mints.
\begin{itemize}
\item a block member is the model's type former at the built identifier, with the block's
parameter and field counts, and conversely;
\item constructors correspond both directions, at inductive and index;
\item a declared inductive's \code{casesOn} is a model constant whose segmentation is \hl{exactly
the arithmetic \code{CasesInfoAgreesK} reads off};
\item a declared inductive is a member of its own block.
\end{itemize}
\lead{Lean} six fields, one spelled \code{casesOnDecl} (no field may be named \code{casesOn}).
\end{definition}

\begin{definition}[The effectful calls only advance the generator]
\label{def:PrimMonotone}
\uses{def:Erasure-EraseM}
\lean{LeanToLambdaBox.PrimMonotone, LeanToLambdaBox.ForallMatchesLam}
\leanok
\lead{In short} Four effect-only calls (reading the environment, logging, checking an instance,
inferring a type) only advance the fresh-name generator; type inference also reports a
\code{forall}-telescope agreeing, in name and domain, with the subject's \code{lambda}-telescope.
\lead{Lean} \code{ForallMatchesLam} is \hl{vacuously true when the subject is not a lambda, by
design}, asserting nothing about a constant's or application's inferred type.
\end{definition}

\begin{definition}[The specification of the eraser's primitives]
\label{def:ErasureSpec}
\uses{def:LookupAdequate, def:PrimMonotone, def:BlockAdequate, def:Oracle-MetaSound, def:isErasable, def:Erasure-isErasable, def:lean4lean-grounding, asm:ErasureSpec-env_connect, asm:ErasureSpec-lookup_adequate, asm:ErasureSpec-fresh_names, asm:ErasureSpec-oracle_refl, asm:ErasureSpec-oracle_meta, asm:ErasureSpec-decl_adequate, asm:ErasureSpec-prim_monotone, asm:ErasureSpec-block_adequate}
\lean{LeanToLambdaBox.ErasureSpec}
\leanok
\lead{In short} Eight \stAssumed{} clauses about the world outside pure Lean terms: an elaboration
environment, its lean4lean model, the universe scope, the fresh-name generator.
\begin{center}
\begin{tabular}{lp{0.68\linewidth}l}
\toprule
Field & What it asserts & Class \\
\midrule
\code{env\_connect} & a well-formed lean4lean model exists, projecting safely & D \\
\code{lookup\_adequate} & Definition~\ref{def:LookupAdequate} & D \\
\code{fresh\_names} & a fresh variable: unreserved, then reserved; generator only advances,
also kernel-reserved & D \\
\code{oracle\_refl} & a positive verdict reflects a pure kernel run, or the assumed-sound fallback
& D \\
\code{oracle\_meta} & a verdict off-scope is sound outright & D \\
\code{decl\_adequate} & a safe declaration is visible in the model, translated & D \\
\code{prim\_monotone} & Definition~\ref{def:PrimMonotone} & D \\
\code{block\_adequate} & Definition~\ref{def:BlockAdequate} & D \\
\bottomrule
\end{tabular}
\end{center}
\lead{Lean} \hl{the top hypothesis of the capstone, every rung, the erasure bridge, the fragment's
soundness lemmas and the cold-start induction}.
\end{definition}

\begin{remark}
\lead{Caveat} A \emph{definition}: the structure elaborates. Each of its eight fields is a
separate assumption node, trust class~D (Chapter~\ref{chap:trust}).
\begin{itemize}
\item \code{decl\_adequate}'s safety reads \code{.safe} $\le$ \code{ci.safety}; since
\code{.safe} is lean4lean's \emph{top}, this forces safety, admitting no unsafe or partial
declaration, though the inequality reads backwards at first sight;
\item no source-table adequacy hypothesis: a \code{Prop} structure cannot hold a table as data.
\end{itemize}
\end{remark}

\begin{remark}
\lead{Caveat} \code{ErasureSpec.lean} also declares a relational pass interface \code{LBPassR}, in
MetaRocq's \code{optimize\_correct} shape. It has no consumer and its doc comment forward-references
a theorem that does not exist: an aspiration, not a current interface.
\end{remark}

\section{What the eraser is asked for: EraserAsks}
\label{sec:eraser:asks}
\code{EraserAsks} names, once, every fact assumed about its own shipping code.

\begin{definition}[What the eraser is asked for]
\label{def:EraserAsks}
\uses{def:preparePasses, def:Erasure-isErasable, def:isErasable, def:Erasure-prepare_erasure, def:SEval, def:IndInfo, def:Kername, asm:EraserAsks-passes_monotone, asm:EraserAsks-passes_sound, asm:EraserAsks-oracle_false_refl, asm:EraserAsks-kernel_ind_head_true, asm:EraserAsks-block_keys_distinct}
\lean{LeanToLambdaBox.EraserAsks}
\leanok
\lead{In short} Five \stAssumed{} clauses about \hl{this repository's own code} -- the
preparation passes and the oracle -- each an obligation with an owner.
\begin{center}
\begin{tabular}{p{0.25\linewidth}p{0.60\linewidth}l}
\toprule
Field & What it asserts & Class \\
\midrule
\code{passes\_monotone} & each call only advances the generator & C \\
\code{passes\_sound} & each pass preserves evaluation under any application spine & C \\
\code{oracle\_false\_refl} & a negative verdict: the pure kernel run did not answer positively
& C \\
\code{kernel\_ind\_head\_true} & at an inductive-type head the pure run answers positively & C \\
\code{block\_keys\_distinct} & a mutual block's members get distinct \lbox{} kernames & C \\
\bottomrule
\end{tabular}
\end{center}
\lead{Lean} \code{oracle\_informative} is a \emph{derived theorem}
(Theorem~\ref{thm:EraserAsks-oracle_informative}), not a sixth field.
\end{definition}

\begin{remark}
\lead{Caveat} Two fields are documented \emph{refuted in general}: \code{kernel\_ind\_head\_true}
by \textbf{F-DEPTH}, \code{block\_keys\_distinct} by \textbf{F-UNSAFEREC}. Both held on a corpus --
$2{,}940$ of $2{,}940$ type formers, no failure over $200{,}000$-plus declaration-info answers --
but since the bundle binds the capstone and all eight rungs, results taking it are conditionally
vacuous on a counterexample.
\end{remark}

\section{What the bundles prove}
\label{sec:eraser:derived}
What the two bundles are worth: the reductions and the one reroute they buy.

\begin{theorem}[The modelled environment is well formed]
\label{thm:ErasureSpec-envWF}
\uses{def:ErasureSpec, def:lean4lean-grounding}
\lean{LeanToLambdaBox.ErasureSpec.envWF}
\leanok
\lead{Given} the specification bundle;
\lead{Then} the model environment is well formed.
\lead{Lean} the bundle is its only hypothesis; \hl{a genuine reduction, not a separate assumption}
-- premise of essentially every target-side lemma of
Chapters~\ref{chap:erases}--\ref{chap:lowering}.
\lead{Status} \stProved{}, conditional on \stAssumed{} \code{ErasureSpec}.
\end{theorem}
\begin{proof}
\leanok
\uses{asm:ErasureSpec-env_connect}
\lead{Idea} Unpack the model-connection field and apply lean4lean's well-formedness at safety
\code{.safe}. Two lines.
\end{proof}

\begin{proposition}[Declaration adequacy, for the kernel's own lookup]
\label{prop:ErasureSpec-decl_adequate_of_kernelFind}
\uses{def:ErasureSpec, def:lean4lean-grounding}
\lean{LeanToLambdaBox.ErasureSpec.decl_adequate_of_kernelFind}
\leanok
\lead{In short} The \code{decl\_adequate} field's content, proved for the \emph{kernel}
environment's own lookup rather than the elaborator's.
\lead{Lean} differs only in which lookup primitive it quantifies over; shows what blocks the rest
-- agreement between the elaborator's and kernel's lookup, an upstream ask. \hl{Proved but never
applied}: not a discharge of the field.
\end{proposition}
\begin{proof}
\leanok
\uses{asm:ErasureSpec-env_connect, asm:lean-reflection-axioms}
\lead{Idea} The model-connection field yields lean4lean's translated-environment predicate; its
lookup lemma reads off a visible declaration.
\end{proof}

\begin{theorem}[The oracle reroute]
\label{thm:ErasureSpec-oracle_sound_of_run}
\uses{def:ErasureSpec, def:Erasure-isErasable, def:Oracle-MetaSound, def:Erasable, def:lean4lean-grounding}
\lean{LeanToLambdaBox.ErasureSpec.oracle_sound_of_run}
\leanok
\lead{Given} the specification bundle, and a run of the shipping oracle returning positive on a
term, at the ambient scope;
\lead{Then} every lean4lean translation of that term, at a context matching the run's and
kernel-reserved, is \hl{\Erasable{}, hence a legitimate target of the box rule}.
\lead{Lean} phrases the verdict as an equation on a concrete \code{MetaM} run.
\lead{Status} \stProved{}, \stInherited{} through lean4lean, conditional on \stAssumed{}
\code{ErasureSpec}.
\end{theorem}
\begin{proof}
\leanok
\uses{thm:Oracle-kernel_isErasable_sound, def:Oracle-MetaSound, asm:ErasureSpec-env_connect, asm:ErasureSpec-oracle_refl, asm:lean4lean-trust, asm:lean-reflection-axioms}
\lead{Idea} Unpack the model connection; \code{oracle\_refl} offers two disjuncts.
\begin{enumerate}
\item Kernel disjunct: apply Theorem~\ref{thm:Oracle-kernel_isErasable_sound}.
\item Fallback disjunct: the conclusion \emph{is} \code{Oracle.MetaSound} at hand.
\end{enumerate}
\end{proof}

\begin{remark}
\lead{Why} The one place the relevance decision stops being assumed and becomes a consequence of a
verified checker. Still assumed: the \emph{reflection} step and the fallback arm.
\begin{itemize}
\item footprint (\code{test/ledger.expected}): 33 axioms -- \code{propext}, \code{sorryAx},
\code{Classical.choice}, \code{Quot.sound}, plus 29 non-standard lean4lean modelling axioms;
\item on record: revert the reroute (forty lines), keep the oracle sound by assumption instead.
\end{itemize}
\end{remark}

\begin{theorem}[The type-former exclusion, derived]
\label{thm:EraserAsks-oracle_informative}
\uses{def:EraserAsks, def:Erasure-isErasable, def:IndInfo, def:lean4lean-grounding}
\lean{LeanToLambdaBox.EraserAsks.oracle_informative}
\leanok
\lead{Given} the \code{EraserAsks} bundle, and the shipping oracle answers negatively at the
ambient scope;
\lead{Then} the head of that term's spine is \hl{not an inductive type name}.
\lead{Lean} concludes for every constant head and instantiation at once, under the run-level
premises of Theorem~\ref{thm:ErasureSpec-oracle_sound_of_run}.
\lead{Status} \stProved{}, conditional on \stAssumed{} \code{EraserAsks}.
\end{theorem}
\begin{proof}
\leanok
\uses{asm:EraserAsks-oracle_false_refl, asm:EraserAsks-kernel_ind_head_true}
\lead{Idea} The bundle's two oracle clauses contradict at an inductive head: one says the pure
kernel run answers positively there, the other that it did not.
\end{proof}

\begin{remark}
\lead{Caveat} A theorem, not a sixth field -- the template for retiring fields: derive rather than
assume. Chapter~\ref{chap:refinement} consumes it, but it is exactly as strong as
\code{kernel\_ind\_head\_true}, refuted in general, and no stronger.
\end{remark}

\begin{lemma}[An inductive-type-name spine is erasable]
\label{lem:erasable_indSpine}
\uses{def:IndInfo, def:Erasable, def:lean4lean-grounding}
\lean{LeanToLambdaBox.erasable_indSpine}
\leanok
\lead{In short} If the model knows a constant as an inductive type former, every application of it
to translating arguments is \Erasable{} -- the box rule \emph{does} cover applied type
constructors.
\lead{Lean} takes the environment and the modelled context well formed and quantifies over all
argument lists at once; \hl{recorded where the code puts it, no consumer anywhere}.
\end{lemma}
\begin{proof}
\leanok
\uses{lem:Erases-indInfo_erasable, thm:Erasable-app, asm:lean4lean-trust}
\lead{Idea} Strong induction on spine length: the head's type is an arity, so erasable, and
erasability carries one argument at a time.
\end{proof}

\section{Reading a constant three ways}
\label{sec:eraser:origin}
Lean has one constant node where Rocq has three; recovered here from the kernel model.

\begin{definition}[Recognising an eliminator by its name]
\label{def:isCasesOnName}
\lean{LeanToLambdaBox.isCasesOnName, LeanToLambdaBox.lastComponent}
\leanok
\lead{In short} A constant is an eliminator exactly when its last name component is
\code{casesOn}.
\lead{Lean} \hl{a deliberately syntactic criterion} -- no kernel ``is this a recursor'' flag is
consulted -- read by Chapter~\ref{chap:source}'s eliminator predicate, the fragment checker, and
two \code{ErasureSpec} clauses.
\end{definition}

\begin{lemma}[Every declared constant is a constructor, a type name or a plain constant]
\label{lem:consts_classified}
\uses{def:UpstreamAsks, def:CtorOf, def:IndInfo, def:ConstOrigin, def:lean4lean-grounding}
\lean{LeanToLambdaBox.consts_classified, LeanToLambdaBox.constOrigin_not_ctorOf, LeanToLambdaBox.constOrigin_not_indInfo, LeanToLambdaBox.constOrigin_of_constants}
\leanok
\lead{In short} Lean has one constant node where Rocq has three, so \hl{the three readings}
(constructor, inductive type name, plain constant) must be recovered; under the upstream ask every
declared constant admits exactly one.
\lead{Lean} \code{consts\_classified}'s well-formedness hypothesis is unused (a leading
underscore); call sites supply it for nothing.
\end{lemma}
\begin{proof}
\leanok
\uses{asm:UpstreamAsks-constsOrigin}
\lead{Idea} Each statement projects the ask's classification conjunct; a pin bump discharges all
at once, with no statement changing shape.
\end{proof}

\begin{lemma}[The coordinates of a name are unique]
\label{lem:IndInfo-inj}
\uses{def:UpstreamAsks, def:IndInfo, def:IndArity, def:CtorOf, def:IndDeclOf, def:CasesOnShape}
\lean{LeanToLambdaBox.IndInfo.inj, LeanToLambdaBox.CtorOf.inj, LeanToLambdaBox.CasesOnShape.inj, LeanToLambdaBox.IndArity.inj, LeanToLambdaBox.IndInfo.indDeclOf, LeanToLambdaBox.indBlock_uniq, LeanToLambdaBox.indBlockBelow_type_uniq, LeanToLambdaBox.addConst_foldlM_name_inj}
\leanok
\lead{In short} \lbox{} coordinates are functions of the source name: a type former's block is
unique (kername, index, counts determined), a constructor belongs to one inductive at one index,
an eliminator's segmentation is unique, and a block names each former once.
\lead{Lean} eight statements, all taking the upstream ask. \hl{Together they identify two
independent witnesses}, so two erasure derivations agree on constructor blocks.
\end{lemma}
\begin{proof}
\leanok
\uses{asm:UpstreamAsks-constsOrigin}
\lead{Idea} The injectivity statements project the ask directly; the block lemmas are the
workhorses, since a successful extension fold has injective names.
\end{proof}

\begin{theorem}[A value of a relevant inductive type is not erasable]
\label{thm:not_erasable_of_informative}
\uses{def:Erasable, def:InformativeInd, def:IndDeclOf, def:UpstreamAsks, def:IndBlockBelow, def:lean4lean-grounding}
\lean{LeanToLambdaBox.not_erasable_of_informative, LeanToLambdaBox.indSpine_not_prop, LeanToLambdaBox.peel_ne_prop, LeanToLambdaBox.indSpine_ne_forallE}
\leanok
\lead{Given} a type former is inductively declared and its result sort never evaluates to
\code{Sort 0} (\emph{informative});
\lead{Then} no term typed at a spine headed by it is erasable: \hl{neither a proof nor a type
former}, so the box rule cannot reach it.
\lead{Lean} takes the environment's well-formedness, the upstream ask, and a context of types.
\lead{Status} \stProved{}, \stInherited{} through lean4lean, conditional on \stAssumed{}
\code{UpstreamAsks}.
\end{theorem}
\begin{proof}
\leanok
\uses{lem:IndInfo-inj, asm:UpstreamAsks-constsOrigin, asm:UpstreamAsks-mkAppsInv, asm:UpstreamAsks-constArityInv, asm:lean4lean-trust}
\lead{Idea} \Erasable{} offers two disjuncts, both refuted.
\begin{enumerate}
\item Proof: would make the spine a proposition, but a never-zero result sort can't peel to
\code{Sort 0}.
\item Type-former: would make the spine defeq to a sort or product, refuted by the declaration
data.
\end{enumerate}
\end{proof}

\begin{remark}
\lead{Why} The non-erasability side condition of the whole development: what lets a first-order
answer survive un-boxed (Chapter~\ref{chap:correctness}). [S, Sec.~7.3]'s erasure relation
excludes the box rule by such a side condition; this is its model-side negation.
\end{remark}

\begin{lemma}[Peeling the typing of a translated spine]
\label{lem:trExprS_spine_peel}
\uses{def:IndBlockBelow, def:lean4lean-grounding}
\lean{LeanToLambdaBox.trExprS_spine_peel}
\leanok
\lead{In short} If a head and a whole application both translate and the head is typed, translated
arguments of the right length exist, and the application's type is reached by peeling the head's
type through them.
\lead{Lean} \hl{takes no upstream ask} -- contrast Theorem~\ref{thm:not_erasable_of_informative}.
\end{lemma}
\begin{proof}
\leanok
\uses{asm:lean4lean-trust}
\lead{Idea} Induction along the spine: lean4lean's translation relation already carries each
step's typing, and unique typing collapses the head's two typings.
\end{proof}

\begin{theorem}[The discriminant of a saturated eliminator has its own inductive type]
\label{thm:elim_major}
\uses{def:CasesOnShape, def:lean4lean-grounding}
\lean{LeanToLambdaBox.elim_major, LeanToLambdaBox.peel_major}
\leanok
\lead{Given} a closed, well-typed application of an inductive's eliminator, segmented as its own
block fixes;
\lead{Then} the discriminant's translation is \hl{typed at a spine headed by that inductive}.
\lead{Lean} at the empty modelled local context; the one kernel corollary of its file taking no
upstream ask.
\lead{Status} \stProved{}, \stInherited{} through lean4lean.
\end{theorem}
\begin{proof}
\leanok
\uses{lem:trExprS_spine_peel, asm:lean4lean-trust}
\lead{Idea} The shape predicate places the major premise at a known position;
Lemma~\ref{lem:trExprS_spine_peel} peels the spine's typing onto the declared type.
\end{proof}

\begin{remark}
\lead{Why} The typing premise of Chapter~\ref{chap:correctness}'s iota-arm simulation: it lets the
scrutinee land on the right constructor, so the \lbox{} case node picks the matching branch. [S,
Sec.~7.4]'s case step.
\end{remark}

\begin{definition}[A product telescope over an inductive spine]
\label{def:PiSpine}
\lean{LeanToLambdaBox.PiSpine}
\leanok
\lead{In short} \code{PiSpine}: a type is a telescope of exactly $n$ products over an application
headed by a given inductive -- \hl{the shape a constructor's declared type has}.
\lead{Lean} defined by recursion on the number of products.
\end{definition}

\begin{lemma}[Peeling a constructor's type exhausts its telescope]
\label{lem:peel_piSpine}
\uses{def:PiSpine, def:IndBlockBelow, def:UpstreamAsks, def:IndDeclOf, def:lean4lean-grounding}
\lean{LeanToLambdaBox.peel_piSpine, LeanToLambdaBox.peel_piSpine_head}
\leanok
\lead{In short} Peeling a type defeq to a product telescope over an inductive spine, at arguments
whose result is again an inductive spine, \hl{consumes the telescope exactly}; peeling to the end
reaches the telescope's result spine.
\lead{Lean} both take the declaredness of both inductives and the upstream ask.
\end{lemma}
\begin{proof}
\leanok
\uses{thm:not_erasable_of_informative, asm:UpstreamAsks-constArityInv, asm:lean4lean-trust}
\lead{Idea} A short spine would end at a product; a long one would need a spine to be a product.
Both refuted by the ask that an inductive spine is defeq to no product type.
\end{proof}

\begin{theorem}[A well-typed constructor spine is saturated]
\label{thm:ctor_saturated}
\uses{def:CtorOf, def:IndInfo, def:IndArity, def:UpstreamAsks, def:PiSpine, def:lean4lean-grounding}
\lean{LeanToLambdaBox.ctor_saturated, LeanToLambdaBox.CtorOf.ctorResult_at}
\leanok
\lead{Given} a closed application of the $k$-th constructor of an inductive, translation typed at
a spine of it;
\lead{Then} it has \hl{exactly as many arguments as the block's parameters plus that constructor's
fields}.
\lead{Lean} at the empty context, taking the constructor's coordinates and arity data explicitly.
\lead{Status} \stProved{}, \stInherited{} through lean4lean, conditional on \stAssumed{}
\code{UpstreamAsks}.
\end{theorem}
\begin{proof}
\leanok
\uses{lem:peel_piSpine, lem:trExprS_spine_peel, lem:IndInfo-inj, asm:UpstreamAsks-constsOrigin, asm:lean4lean-trust}
\lead{Idea} The constructor's type is a telescope of that length over the inductive's spine.
\begin{enumerate}
\item Identify the constructor's block with the arity's block by uniqueness.
\item Peel the translated spine's typing onto it (Lemma~\ref{lem:trExprS_spine_peel}).
\item Force the length (Lemma~\ref{lem:peel_piSpine}).
\end{enumerate}
\end{proof}

\begin{remark}
\lead{Deviation} Would license \lbox{}'s \emph{block} form, MetaRocq's
\code{constructors\_as\_blocks} counterpart -- yet this eraser emits applied form anyway
(Definition~\ref{def:Erasure-visitExpr}). A reader expecting block form must be told otherwise.
\end{remark}

\begin{theorem}[The eliminator's block is the inductive's block]
\label{thm:CasesOnShape-agree}
\uses{def:CasesOnShape, def:UpstreamAsks, def:IndArity, def:lean4lean-grounding}
\lean{LeanToLambdaBox.CasesOnShape.agree, LeanToLambdaBox.CasesOnShape.indBlockBelow, LeanToLambdaBox.IndArity.indBlockBelow}
\leanok
\lead{Given} a constant is an inductive's eliminator at a segmentation, and the inductive has
source-side arity data;
\lead{Then} the pre-discriminant arguments decompose as parameters, one motive, the indices;
minor premises equal the constructor count.
\lead{Lean} \hl{lets the source-side iota rule and the \lbox{} case node be about the same block}
-- the pivot of Chapter~\ref{chap:correctness}'s iota arm; its measured footprint has no
\code{sorryAx}.
\lead{Status} \stProved{}, conditional on \stAssumed{} \code{UpstreamAsks}.
\end{theorem}
\begin{proof}
\leanok
\uses{lem:IndInfo-inj, asm:UpstreamAsks-constsOrigin}
\lead{Idea} Each hypothesis exhibits a well-formed declaration list; block uniqueness identifies
the two, and the arithmetic is read off the same block twice.
\end{proof}

\begin{lemma}[A tabled constant is a constant of the model]
\label{lem:constants_of_tabled}
\uses{def:ErasureSpec, def:Witness-SourceTableAdequate, def:TableSafe, def:lean4lean-grounding}
\lean{LeanToLambdaBox.constants_of_tabled}
\leanok
\lead{In short} A constant the reified compiler table gives a body for is a declared constant of
the model.
\lead{Lean} takes the bundle, the table's adequacy and safety hypothesis, concluding \hl{only
existence of the model constant, not its shape}.
\end{lemma}
\begin{proof}
\leanok
\uses{asm:ErasureSpec-decl_adequate}
\lead{Idea} Table adequacy pins the constant in the environment; safety gives its level;
\code{decl\_adequate} translates it into the model.
\end{proof}

\begin{remark}
\lead{Caveat} One of two halves Chapter~\ref{chap:lowering} needs at a tabled constant; the other
is \code{constOrigin\_of\_constants} (Lemma~\ref{lem:consts_classified}). Their join -- neither
constructor nor type name -- is an upstream ask, \hl{not proved here}, standing as a premise until
the pin bump lands.
\end{remark}

\section{Deviations from the paper, and caveats}
\label{sec:eraser:caveats}

\begin{center}
\begin{tabular}{lp{0.24\linewidth}p{0.50\linewidth}}
\toprule
Tag & Site & Symptom \\
\midrule
F-DEPTH & arity-check fuel & exhausts behind \code{@[irreducible]}; falls back, negative \\
F-UNSAFEREC & \code{remove\_unsafe\_rec} & not injective: can collapse two kernames of a legal
block \\
F-PROP & \code{register\_inductive} & \code{propositional}, \code{kelim}, \code{finite} stay
default; match-on-box never fires \\
F-SPARSE & \code{visitCases} & sparse eliminator declared unreachable; writes a wrong \code{.ast} \\
F-ETA2 & \code{visitConstructor} eta loops & re-erases the argument prefix per missing argument,
for nothing \\
F-ETA & \code{visitMutual} & bodies ship bare, unapplied; expanded-program predicate fails on all
five benchmark programs \\
F-QUOT, F-EQREC & \code{visitMutual} axiom arm & \code{Quot}/\code{Eq.rec} constants land as
body-less axioms \\
\bottomrule
\end{tabular}
\end{center}

\begin{itemize}
\item \lead{Deviation} Constructors ship applied, not blocked -- opposite agda2lambox and
\code{constructors\_as\_blocks} -- pinned off throughout.
\item \lead{Deviation} Fixpoint bodies ship bare and unapplied (F-ETA); the capstone concludes
peregrine well-formedness, not MetaRocq's expanded-program form.
\item \lead{Deviation} The preparation pass has no MetaCoq counterpart; \code{ConfigPinned} pins
its \code{@[csimp]} walk off.
\item \lead{Deviation} The oracle reroute re-implements MetaRocq's test on lean4lean's checker, at
a 33-name axiom footprint.
\item \lead{Caveat} Two \code{EraserAsks} fields are refuted in general; results using the bundle
are conditionally vacuous on a counterexample.
\item \lead{Caveat} \code{panic!} yields $\Box$ and exits successfully, so F-SPARSE and others
surface as a wrong \code{.ast}, not a failure.
\end{itemize}
